\chapter{2010年天津卷（文）}

\section{单选题}

\begin{problem}
1 是虚数单位，复数 \(\frac{3 + \i}{1 - \i} =\)（\quad）

\choices
    {\(1 + 2\i\)}
    {\(2 + 4\i\)}
    {\(-1\) -2i}
    {\(2 - \i\)}
\end{problem}

\begin{problem}
设变量 \(x, y\) 满足约束条件 \(\left\{ \begin{array}{l} x + y \leqslant 3, \\ x - y \geqslant -1, \\ y \geqslant 1, \end{array} \right.\) 则目标函数 \(z = 4x + 2y\) 的最大 值为（\quad）

\choices
    {12}
    {10}
    {8}
    {2}
\end{problem}

\begin{problem}
阅读下面的程序框图，若输出 \(s\) 的值为 \(-7\)，则解析框内可填写（\quad）

\choices
    {\(-1\)}
    {0}
    {1}
    {3}
\end{problem}

\begin{problem}
函数 \( f(x) = \e^{x} + x - 2 \) 的零点所在的一个区间是（\quad）

\choices
    {\((-2, -1)\)}
    {\((-1,0)\)}
    {\((0,1)\)}
    {(1,2)}
\end{problem}

\begin{problem}
下列命题中，真命题是（\quad）

\choices
    {\(3m \in \R\)，使函数 \(f(x) = x^{2} + mx (x \in \R)\) 是偶函数}
    {\(3m \in \R\)，使函数 \(f(x) = x^{2} + mx (x \in \R)\) 是奇函数}
    {\(\forall m \in \R\)，使函数 \(f(x) = x^2 + mx (x \in \R)\) 是偶函数}
    {\(\forall m \in \R\)，使函数 \(f(x) = x^2 + mx (x \in \R)\) 是奇函数}
\end{problem}

\begin{problem}
设 \(a = \log_64, b = (\log_63)^2, c = \log_45\)，则（\quad）

\choices
    {\(a < c < b\)}
    {\(b < c < a\)}
    {\(a < b < c\)}
    {\(b < a < c\)}
\end{problem}

\begin{problem}
设集合 \(A = \{x \mid |x - a| < 1, x \in \R\}\)，\(B = \{x \mid 1 < x < 5, x \in \R\}\). 若 \(A \cap B = \varnothing\)，则实数 \(a\) 的取值范围是（\quad）

\choices
    {\(\{a\mid 0\leqslant a\leqslant 6\}\)}
    {\(\{a\mid a\leqslant 2\) 或 \(a\geqslant 4\}\)}
    {\(\{a\mid a\leqslant 0\) 或 \(a\geqslant 6\}\)}
    {\(\{a\mid 2\leqslant a\leqslant 4\}\)}
\end{problem}

\begin{problem}
如图是函数 \( y = A \sin (\omega x + \varphi) (x \in \R) \) 在区间 \( \left[-\frac{\pi}{6}, \frac{5\pi}{6}\right] \) 上的图象，为了得到这个函数的图象，只要将 \( y = \sin x (x \in \R) \) 的图象上所有的点（\quad）

\choices
    {向左平移 \(\frac{\pi}{3}\) 个单位长度，再把所得各点的横坐标缩短到原来的 \(\frac{1}{2}\) 倍，}
    {向左平移 \(\frac{\pi}{3}\) 个单位长度，再把所得各点的横坐标伸长到原来的 2 倍，纵坐标不变}
    {向左平移 \(\frac{\pi}{6}\) 个单位长度，再把所得各点的横坐标缩短到原来的 \(\frac{1}{2}\) 倍，纵坐标不变}
    {向左平移 \(\frac{\pi}{6}\) 个单位长度，再把所得各点的横坐标伸长到原来的 2 倍，纵坐标不变}
\end{problem}

\begin{problem}
如图，在 \(\triangle ABC\) 中，\(AD \perp AB, \overrightarrow{BC} = \sqrt{3}\overrightarrow{BD}, |\overrightarrow{AD}| = 1\)，则 \(\overrightarrow{AC} \cdot \overrightarrow{AD} =\)（\quad）

\choices
    {\(2\sqrt{3}\)}
    {\(\frac{\sqrt{3}}{2}\)}
    {\(\frac{\sqrt{3}}{3}\)}
    {\(\sqrt{3}\)}
\end{problem}

\begin{problem}
设函数 \( g(x) = x^{2} - 2 (x \in \R) \)，\( f(x) = \begin{cases} g(x) + x + 4, & x < g(x), \\ g(x) - x, & x \geqslant g(x), \end{cases} \) 则 \( f(x) \) 的值域是（\quad）

\choices
    {\(\left[-\frac{9}{4}, 0\right] \cup (1, +\infty)\)}
    {\([0, +\infty)\)}
    {\(\left[-\frac{9}{4}, +\infty\right)\)}
    {\(\left[-\frac{9}{4}, 0\right] \cup (2, +\infty)\)}
\end{problem}

\section{填空题}

\begin{problem}
如图，四边形 \(ABCD\) 是圆 \(O\) 的内接四边形，延长 \(AB\) 和 \(DC\) 相交于点 \(P\). 若 \(PB = 1, PD = 3\)，则 \(\frac{AC}{DC}\) 的值为 \fillinblank{} \fillinblank{}.
\end{problem}

\begin{problem}
一个几何体的三视图如图所示，则这个几何体的体积为 \fillinblank{}.

正视图

侧视图

俯视图
\end{problem}

\begin{problem}
已知双曲线 \(\frac{x^2}{a^2} - \frac{y^2}{b^2} = 1 (a > 0, b > 0)\) 的一条渐近线方程是 \(y = \sqrt{3} x\)，它的一个焦点与抛物线 \(y^2 = 16x\) 的焦点相同. 则双曲线的方程为 \fillinblank{} \fillinblank{}.
\end{problem}

\begin{problem}
已知圆 \(C\) 的圆心是直线 \(x - y + 1 = 0\) 与 \(x\) 轴的交点，且圆 \(C\) 与直线 \(x + y + 3 = 0\) 相切，则圆 \(C\) 的方程为 \fillinblank{} \fillinblank{}.
\end{problem}

\begin{problem}
设 \(\{\alpha_{n}\}\) 是等比数列，公比 \(q = \sqrt{2}\). \(S_{n}\) 为 \(\{\alpha_{n}\}\) 的前 \(n\) 项和. 记 \(T_{n} = \frac{17S_{n} - S_{n}}{a_{n + 1}}, n \in \mathbf{N}^{*}\). 设 \(T_{100}\) 为数列 \(\{T_{n}\}\) 的最大项，则 \(n_{99} = \fillinblank{}\).
\end{problem}

\begin{problem}
设函数 \( f(x) = x - \frac{1}{x} \)，对任意 \( x \in [1, +\infty) \)，\( f(mx) + mf(x) < 0 \) 恒成立，则实数 \( m \) 的取值范围是 \fillinblank{} \fillinblank{}.
\end{problem}

\section{解答题}

\begin{problem}
在 \(\triangle ABC\) 中: \(\frac{AC}{AB} = \frac{\cos B}{\cos C}\). (1) 证明： \( B = C = \frac{1}{3} \) (2) 若 \(\cos A = -\frac{1}{3}\)，求 \(\sin \left(4B + \frac{\pi}{3}\right)\) 的值.
\end{problem}

\begin{problem}
有编号为 \(A_{1}, A_{2}, \cdots, A_{10}\) 的 10 个零件，测量其直径 (单位: cm)，得到下面数据: 编号 \( A_1 \) \( A_2 \) \( A_3 \) \( A_4 \) \( A_5 \) \( A_6 \) \( A_7 \) \( A_8 \) \( A_9 \) \( A_{10} \) 直径 1.51 1.49 1.49 1.51 1.49 1.51 1.47 1.46 1.53 1.47 其中直径在区间 [1.48, 1.52] 内的零件为一等品. (1) 从上述 10 个零件中，随机抽取一个，求这个零件为一等品的概率; (2) 从一等品零件中，随机抽取 2 个. \circled{1} 用零件的编号列出所有可能的抽取结果； \circled{2} 求这 2 个零件直径相等的概率.
\end{problem}

\begin{problem}
已知函数 \( f(x) = ax^3 - \frac{3}{2} x^2 + 1 (x \in \R) \)，其中 \( a > 0 \). (1) 若 \(a = 1\)，求曲线 \(y = f(x)\) 在点 (2, \(f(2)\)) 处的切线方程; (2) 若在区间 \(\left[\frac{1}{2}, \frac{1}{2}\right]\) 上，\(f(x) > 0\) 恒成立，求 \(a\) 的取值范围.
\end{problem}

\begin{problem}
在数列 \(\{a_{n}\}\) 中，\(a_{1} = 0\)，且对任意 \(k \in \mathbf{N}^{*}\)，\(a_{2k-1}, a_{2k}, a_{2k+1}\) 成等差数列，其公差为 \(2k\). (1) 证明 \(a_{k}, a_{2k}, a_{3k}\) 成等比数列； (2) 求数列 \(\{a_{n}\}\) 的通项公式; (3) 记 \(T_{n} = \frac{2^{n}}{a_{2}} + \frac{3^{n}}{a_{3}} + \dots + \frac{n^{2}}{a_{m}}\)，证明 \(\frac{3}{2} < 2n - T_{n} \leqslant 2 (n \geqslant 2)\).
\end{problem}

\begin{problem}
如图，在五面体 \(ABCDEF\) 中，四边形 \(ADEF\) 是正方形，\(FA \perp\) 平面 \(ABCD\)，\(BC \nmid AD\)，\(CD = 1\)，\(AD = 2\sqrt{2}\)，\(\angle BAD = \angle CDA = 45^\circ\). (1) 求异面直线 \(CE\) 与 \(AF\) 所成角的余弦值; (2) 证明 \(CD \perp\) 平面 \(ABF\); (3) 求二面角 \(B - EF - A\) 的正切值.
\end{problem}

\begin{problem}
已知椭圆 \(\frac{x^2}{a^2} + \frac{y^2}{b^2} = 1 (a > b > 0)\) 的离心率 \(e = \frac{\sqrt{3}}{2}\)，连接椭圆的四个顶点得到的菱形的面积为 4. (1) 求椭圆的方程; (2) 设直线 \(l\) 与椭圆相交于不同的两点 \(A, B\)，已知点 \(A\) 的坐标为 \((-a, 0)\). \circled{1} 若 \( |AB| = \frac{4\sqrt{3}}{5} \)，求直线 \( l \) 的倾斜角； \circled{2} 若点 \(Q(0, y_0)\) 在线段 \(AB\) 的垂直平分线上，且 \(\overrightarrow{QA} \cdot \overrightarrow{QB} = 4\). 求 \(y_0\) 的值.
\end{problem}

